% TOP_PROBLEM: {"id":30007737,"problem_number":"LOCAL-30007737","title":"Causes of persistently low fertility","category_id":21,"category":{"id":21,"name":"economics","display_name":"Economics","description":"Open economic theory statements, published research agendas, and proposed scoped specifications; question provenance and historical priority are qualified per record.","slug":"economics","order_index":21,"created_at":"2026-10-08T00:00:00Z"},"status":"provenance_pending","external_url":null,"legacy_ids":[],"published":false,"statement":"Identify or bound how prices, perceived norms, and work, partnership, and leisure opportunities contribute to completed-fertility changes across specified cohorts.","economics":{"original_id":226,"field":"Education, families and demography","kind":"identification","formalization_basis":"proposed_specification","source_status":"not_verified","priority_status":"not_established","priority_note":"No readable primary passage explicitly stating this precise open question was verified. A supporting paper, abstract, or topic citation is insufficient to establish origin or unresolved status.","earliest_verified_reference":null,"current_reference":{"authors":["Melissa S. Kearney","Phillip B. Levine"],"title":"Why Is Fertility So Low in High-Income Countries?","year":2026,"url":"https://www.aeaweb.org/articles?id=10.1257/jel.20261786","doi":"10.1257/jel.20261786","locator":"Abstract","evidence_summary":"The published review discusses explanations for falling fertility; the proposed counterfactual decomposition and completed-fertility policy contrast are our specifications."},"formalization_note":"Proposed scoped research specification. The original catalogue prompt was a synthesis, and no canonical conjecture or verified original formulation is claimed. The supporting source, when retained, establishes context or existing results rather than this exact open question."},"statusQualification":"Provenance pending: an explicit published statement of this exact open question was not verified. This is a catalogue-authored candidate specification, not a certified open conjecture. Proposed scoped research specification. The original catalogue prompt was a synthesis, and no canonical conjecture or verified original formulation is claimed. The supporting source, when retained, establishes context or existing results rather than this exact open question.","statusReviewedAt":"2026-10-08"}
% FIRST_POSTED: 2026-10-08
% ENTRY_KIND: statement_only
% !TeX program = lualatex
\documentclass[11pt]{article}
\usepackage{fontspec}
\usepackage[a4paper,margin=25mm]{geometry}
\usepackage{amsmath,amssymb,mathtools}
\usepackage{xurl}
\usepackage[hidelinks]{hyperref}
\newcommand{\sref}[2]{\hyperlink{source:#1:#2}{[S#2]}}
\setlength{\emergencystretch}{3em}
\begin{document}
% problemId:30007737
\section*{Causes of persistently low fertility}\label{30007737}
\subsection{Definitions and mathematical statement}
Consider women in birth cohorts 1970–1990 in a prespecified set of high-income countries, observed through age forty-five. Let \(F_i(p,n,\ell)\) be completed live births under a jointly specified history of child-related prices \(p\), perceived parenting and gender norms \(n\), and work, partnership, and leisure opportunities \(\ell\). Baseline histories \((p_0,n_0,\ell_0)\) are taken from an earlier cohort; observed later histories are \((p_1,n_1,\ell_1)\). Define the total counterfactual cohort contrast
\[\Delta=E[F_i(p_1,n_1,\ell_1)-F_i(p_0,n_0,\ell_0)].\]
Specify how age-specific partnership formation, desired children, infertility, and postponement enter fertility decisions. Estimate or bound contributions of prices, norms, and opportunities using a prespecified decomposition order or an average across all orders, since interactions prevent a unique additive attribution. Norms must have observable indicators and an explicit formation process; unexplained residuals cannot simply be labeled preference change. State policy and cohort variation supporting causal identification, and distinguish period birth rates from completed fertility. The cited review motivates these competing explanations and future research, but the precise decomposition above was not verified in a full-text open-question passage. This remains a proposed empirical formalization.

\par\textbf{Formulation and scope.} Proposed scoped research specification. The original catalogue prompt was a synthesis, and no canonical conjecture or verified original formulation is claimed. The supporting source, when retained, establishes context or existing results rather than this exact open question.

\par\textbf{Open-question provenance.} An explicit published open-question statement was not verified. Sources below are supporting literature and must not be treated as a first listing.

\par\textbf{Historical priority.} No readable primary passage explicitly stating this precise open question was verified. A supporting paper, abstract, or topic citation is insufficient to establish origin or unresolved status.

\subsection{Short English statement}
Identify or bound how prices, perceived norms, and work, partnership, and leisure opportunities contribute to completed-fertility changes across specified cohorts.

\subsection{Sources}
\begin{itemize}\raggedright
\item[\textnormal{[S1]}]\hypertarget{source:30007737:1}{} Melissa S. Kearney, Phillip B. Levine. 2026. Why Is Fertility So Low in High-Income Countries?. Abstract.
\url{https://www.aeaweb.org/articles?id=10.1257/jel.20261786}.
DOI: 10.1257/jel.20261786.
{\small\textit{Supporting or current literature; see evidence qualification. The published review discusses explanations for falling fertility; the proposed counterfactual decomposition and completed-fertility policy contrast are our specifications.}}
\end{itemize}
\end{document}
